\chapter{Algebra: Equations and Inequalities}\label{ch:g10:algebra}

Algebra is the art of computing with letters. This chapter reviews and
sharpens the two basic moves --- expanding and factoring --- and uses them
to solve equations and inequalities systematically, with the sign table as
the central tool.

\section{Expanding and factoring}

\begin{definition}[Expanding, factoring]\label{def:g10:algebra:expand}
\emph{Expanding}\index{expanding} an expression means transforming
products into sums; \emph{factoring}\index{factoring} means the reverse,
transforming a sum into a product. The basic rule is distributivity:
\[
k(a + b) = ka + kb,
\qquad
(a + b)(c + d) = ac + ad + bc + bd .
\]
\end{definition}

\begin{example}\label{ex:g10:algebra:expand}
Expand $(2x + 3)(x - 5)$ step by step:
\begin{align*}
(2x + 3)(x - 5)
&= 2x \times x + 2x \times (-5) + 3 \times x + 3 \times (-5) \\
&= 2x^2 - 10x + 3x - 15 \\
&= 2x^2 - 7x - 15 .
\end{align*}
\end{example}

\begin{theorem}[Remarkable identities]\label{thm:g10:algebra:identities}
For all real numbers $a$ and $b$:
\[
(a+b)^2 = a^2 + 2ab + b^2,
\qquad
(a-b)^2 = a^2 - 2ab + b^2,
\qquad
(a+b)(a-b) = a^2 - b^2 .
\]
\end{theorem}

\begin{proof}
Each one is a direct expansion. For the first:
\[
(a+b)^2 = (a+b)(a+b) = a^2 + ab + ba + b^2 = a^2 + 2ab + b^2 .
\]
The second follows by replacing $b$ with $-b$: $(a-b)^2 = a^2 + 2a(-b) +
(-b)^2 = a^2 - 2ab + b^2$. For the third:
$(a+b)(a-b) = a^2 - ab + ba - b^2 = a^2 - b^2$.
\end{proof}

\begin{omfigure}
\begin{tikzpicture}[scale=0.62]
  \draw[thick] (0,0) rectangle (5,5);
  \draw[thick] (3.4,0) -- (3.4,5);
  \draw[thick] (0,3.4) -- (5,3.4);
  \fill[omDef!12] (0,0) rectangle (3.4,3.4);
  \fill[omThm!12] (3.4,3.4) rectangle (5,5);
  \fill[omProp!12] (3.4,0) rectangle (5,3.4);
  \fill[omProp!12] (0,3.4) rectangle (3.4,5);
  \node[font=\small] at (1.7,1.7) {$a^2$};
  \node[font=\small] at (4.2,4.2) {$b^2$};
  \node[font=\small] at (4.2,1.7) {$ab$};
  \node[font=\small] at (1.7,4.2) {$ab$};
  \node[below, font=\small] at (1.7,0) {$a$};
  \node[below, font=\small] at (4.2,0) {$b$};
  \node[left, font=\small] at (0,1.7) {$a$};
  \node[left, font=\small] at (0,4.2) {$b$};
\end{tikzpicture}

{\small Why $(a+b)^2 = a^2 + 2ab + b^2$: a square of side $a + b$ splits
into a square of side $a$, a square of side $b$, and two $a \times b$
rectangles.}
\end{omfigure}

\begin{example}\label{ex:g10:algebra:identities}
Used forward (expanding) and backward (factoring):
\begin{align*}
(3x + 2)^2 &= 9x^2 + 12x + 4,
& x^2 - 25 &= (x+5)(x-5), \\
(x - 4)^2 &= x^2 - 8x + 16,
& 4x^2 + 4x + 1 &= (2x + 1)^2 .
\end{align*}
Mental arithmetic benefits too: $101 \times 99 = (100+1)(100-1) =
10000 - 1 = 9999$.
\end{example}

\begin{method}[Factoring an expression]\label{met:g10:algebra:factor}
Try, in this order:
\begin{enumerate}
  \item \emph{Common factor:} spot a factor present in every term and pull
        it out, e.g.\ $6x^2 + 4x = 2x(3x + 2)$. The common factor can be a
        whole bracket: $(x+1)(2x-3) + (x+1)(x+7) = (x+1)(3x+4)$.
  \item \emph{Remarkable identity:} recognize $a^2 - b^2$, or
        $a^2 \pm 2ab + b^2$, e.g.\ $9x^2 - 16 = (3x)^2 - 4^2 =
        (3x+4)(3x-4)$.
  \item Combine both, and check the result by expanding it back.
\end{enumerate}
\end{method}

\begin{example}\label{ex:g10:algebra:factorsteps}
Factor $E = (2x + 1)^2 - (x - 3)^2$. This is a difference of two squares
$a^2 - b^2$ with $a = 2x+1$ and $b = x-3$:
\begin{align*}
E &= \bigl[(2x+1) + (x-3)\bigr] \times \bigl[(2x+1) - (x-3)\bigr] \\
  &= (3x - 2)(2x + 1 - x + 3) \\
  &= (3x - 2)(x + 4).
\end{align*}
Mind the second bracket: subtracting $x - 3$ means subtracting $x$
\emph{and} adding $3$.
\end{example}

\section{Equations}

\begin{definition}[Equation, solution]\label{def:g10:algebra:equation}
An \emph{equation}\index{equation} is an equality containing an unknown
number $x$; a \emph{solution} is a value of $x$ that makes the equality
true. Solving the equation means finding \emph{all} its solutions.
\end{definition}

Two equations are \emph{equivalent} when they have exactly the same
solutions. Adding the same number to both sides, or multiplying both sides
by the same \emph{nonzero} number, produces an equivalent equation.

\begin{example}[First-degree equation]\label{ex:g10:algebra:firstdegree}
Solve $5x - 7 = 2x + 8$, one move at a time:
\begin{align*}
5x - 7 &= 2x + 8 \\
5x - 2x &= 8 + 7 && \text{(add $7 - 2x$ to both sides)} \\
3x &= 15 \\
x &= 5 && \text{(divide both sides by $3$).}
\end{align*}
The unique solution is $5$. Check: $5 \times 5 - 7 = 18$ and
$2 \times 5 + 8 = 18$.
\end{example}

\begin{theorem}[Zero-product rule]\label{thm:g10:algebra:zeroproduct}
A product of real numbers is zero if and only if at least one of its
factors is zero:
\[
A \times B = 0
\quad\Longleftrightarrow\quad
A = 0 \ \text{ or } \ B = 0 .
\]
\end{theorem}

\begin{proof}
If $A = 0$ or $B = 0$, the product is clearly $0$. Conversely, suppose
$AB = 0$ and $A \neq 0$. Dividing both sides by $A$ gives $B = 0$.
\end{proof}

\begin{method}[Solving by factoring]\label{met:g10:algebra:zeroproduct}
To solve an equation whose right-hand side is $0$ and whose left-hand
side can be factored:
\begin{enumerate}
  \item bring everything to the left so the equation reads $E = 0$;
  \item factor $E$ into a product of first-degree factors;
  \item apply the zero-product rule and solve each small equation;
  \item collect all the solutions found.
\end{enumerate}
\end{method}

\begin{example}\label{ex:g10:algebra:zeroproduct}
Solve $(x - 2)(3x + 6) = 0$: either $x - 2 = 0$, giving $x = 2$, or
$3x + 6 = 0$, giving $x = -2$. Solutions: $-2$ and $2$.

Solve $x^2 = 9x$. Do \emph{not} divide by $x$ (it could be zero!); bring
everything to one side and factor:
\[
x^2 - 9x = 0
\quad\Longleftrightarrow\quad
x(x - 9) = 0
\quad\Longleftrightarrow\quad
x = 0 \ \text{ or } \ x = 9 .
\]
\end{example}

\begin{example}[Quotient equation]\label{ex:g10:algebra:quotient}
Solve $\dfrac{x - 1}{x + 2} = 0$. A quotient is zero exactly when its
numerator is zero \emph{and} its denominator is not. Here $x - 1 = 0$
gives $x = 1$, and $1 + 2 = 3 \neq 0$: the only solution is $1$. The value
$x = -2$ is \emph{forbidden} (division by zero) and must always be
excluded first.
\end{example}

\section{Inequalities and sign tables}

\begin{proposition}[Rules for inequalities]\label{prop:g10:algebra:ineqrules}
Let $a \leq b$.
\begin{enumerate}
  \item For any $c$: $a + c \leq b + c$ (adding preserves the order).
  \item If $c > 0$: $ac \leq bc$ (multiplying by a positive number
        preserves the order).
  \item If $c < 0$: $ac \geq bc$ (multiplying by a negative number
        \emph{reverses} the order).
\end{enumerate}
\end{proposition}

\begin{proof}
1. $(b + c) - (a + c) = b - a \geq 0$.
2. $bc - ac = (b-a)c$ is the product of two nonnegative numbers, so
$bc \geq ac$.
3. Now $(b-a)c \leq 0$ as the product of a nonnegative and a negative
number, so $bc \leq ac$.
\end{proof}

\begin{example}\label{ex:g10:algebra:firstineq}
Solve $-2x + 3 > 9$:
\begin{align*}
-2x + 3 &> 9 \\
-2x &> 6 && \text{(subtract 3)} \\
x &< -3 && \text{(divide by $-2 < 0$: reverse the inequality).}
\end{align*}
The solution set is $\intoo{-\infty}{-3}$.
\end{example}

\begin{definition}[Sign table]\label{def:g10:algebra:signtable}
A \emph{sign table}\index{sign table} records, on one row per factor, the
sign ($+$ or $-$) of each factor of an expression as $x$ runs over $\R$,
with a $0$ at each value where the factor vanishes; a final row gives the
sign of the product, using the rule of signs column by column.
\end{definition}

\begin{example}\label{ex:g10:algebra:signtable}
Sign of $P(x) = (x - 2)(3 - x)$. The factor $x - 2$ vanishes at $2$ and is
positive after; the factor $3 - x$ vanishes at $3$ and is positive before:
\begin{center}
\renewcommand{\arraystretch}{1.3}
\begin{tabular}{c|ccccc}
$x$ & $-\infty$ \quad & $2$ & & $3$ & \quad $+\infty$ \\
\hline
$x - 2$ & $-$ & $0$ & $+$ & & $+$ \\
$3 - x$ & $+$ & & $+$ & $0$ & $-$ \\
\hline
$P(x)$ & $-$ & $0$ & $+$ & $0$ & $-$ \\
\end{tabular}
\end{center}
So $P(x) > 0$ exactly on $\intoo{2}{3}$, and $P(x) \leq 0$ on
$\intoc{-\infty}{2} \cup \intco{3}{+\infty}$.
\end{example}

\begin{omfigure}
\begin{tikzpicture}
\begin{axis}[omaxis,
  width=8.2cm, height=5.2cm,
  xmin=-0.6, xmax=5.4, ymin=-2.6, ymax=1.4,
  xtick={2,3}, ytick=\empty]
  \addplot[omDef, thick, domain=0.35:4.65] {(x-2)*(3-x)};
  \fill (axis cs:2,0) circle (1.5pt);
  \fill (axis cs:3,0) circle (1.5pt);
  \node[omDef, font=\small] at (axis cs:2.5,0.65) {$P > 0$};
  \node[omThm, font=\small] at (axis cs:0.9,-1.3) {$P < 0$};
  \node[omThm, font=\small] at (axis cs:4.1,-1.3) {$P < 0$};
\end{axis}
\end{tikzpicture}

{\small The graph of $P(x) = (x-2)(3-x)$ confirms the sign table: positive
between the roots $2$ and $3$, negative outside.}
\end{omfigure}

\begin{method}[Solving an inequality with a sign table]\label{met:g10:algebra:signtable}
To solve an inequality like $E(x) \geq 0$ or $E(x) < 0$:
\begin{enumerate}
  \item bring everything to the left-hand side, so the right-hand side is
        $0$ (never multiply both sides by an expression whose sign is
        unknown);
  \item factor the left-hand side, or put it over a common denominator if
        it is a quotient;
  \item build the sign table with one row per factor (for a quotient, mark
        forbidden values with a double bar in the last row);
  \item read off the intervals where the requested sign holds, checking
        whether each bound is included.
\end{enumerate}
\end{method}

\begin{example}\label{ex:g10:algebra:quotientineq}
Solve $\dfrac{x + 1}{x - 3} \leq 0$. The numerator vanishes at $-1$, the
denominator at $3$ (forbidden value). The quotient is negative when the
two have opposite signs, which happens for $x$ between $-1$ and $3$.
Including the bound where the \emph{numerator} vanishes but excluding the
forbidden value: the solution set is $\intco{-1}{3}$.
\end{example}

\section{Exercises}

\begin{exercise}[$\star$]\label{exo:g10:algebra:1}
Expand and simplify:
\[
3(2x - 5) - 2(x - 7), \qquad
(x + 4)(2x - 1), \qquad
(3x - 2)^2, \qquad
(5 - x)(5 + x).
\]
\end{exercise}

\begin{exercise}[$\star$]\label{exo:g10:algebra:2}
Factor:
\[
x^2 - 49, \qquad
x^2 + 10x + 25, \qquad
7x^2 - 14x, \qquad
16x^2 - 8x + 1 .
\]
\end{exercise}

\begin{exercise}[$\star$]\label{exo:g10:algebra:3}
Solve:
\[
4x + 9 = x - 3, \qquad
\frac{2x - 1}{3} = 5, \qquad
2(x - 3) = 2x + 1 .
\]
(One of them has no solution --- explain why.)
\end{exercise}

\begin{exercise}[$\star$]\label{exo:g10:algebra:4}
Solve using the zero-product rule:
\[
(x + 5)(2x - 8) = 0, \qquad
x^2 - 16 = 0, \qquad
x^2 + 6x + 9 = 0 .
\]
\end{exercise}

\begin{exercise}[$\star$]\label{exo:g10:algebra:5}
Solve the inequalities and write the solution sets as intervals:
\[
3x - 5 \leq 7, \qquad
-4x + 1 > 9, \qquad
2(x - 1) \geq 5x + 4 .
\]
\end{exercise}

\begin{exercise}[$\star\star$]\label{exo:g10:algebra:6}
Factor $E = (x - 1)(x + 3) - (x - 1)(2x - 5)$, then solve $E = 0$.
\end{exercise}

\begin{exercise}[$\star\star$]\label{exo:g10:algebra:7}
Solve $x^2 = 5x$, then $(2x + 1)^2 = (x - 4)^2$. (For the second one,
bring everything to the left and factor the difference of squares.)
\end{exercise}

\begin{exercise}[$\star\star$]\label{exo:g10:algebra:8}
Using a sign table, solve
\[
(x + 2)(4 - x) > 0,
\qquad
\frac{2x - 6}{x + 1} \geq 0 .
\]
\end{exercise}

\begin{exercise}[$\star\star$]\label{exo:g10:algebra:9}
A streaming service costs $9$ per month, plus a one-time signup fee of
$15$; a rival costs $12$ per month with no fee. After how many months does
the first service become the cheaper choice overall? Set up and solve an
inequality.
\end{exercise}

\begin{exercise}[$\star\star$]\label{exo:g10:algebra:10}
Solve $\dfrac{x+3}{x-1} = 2$. (Multiply both sides by $x - 1$ after
excluding the forbidden value, or bring everything to the left over a
common denominator.)
\end{exercise}

\begin{exercise}[$\star\star\star$]\label{exo:g10:algebra:11}
Let $n$ be an integer. Show that $(n + 1)^2 - n^2$ is always odd, and that
the difference between the squares of two consecutive odd numbers is
always a multiple of $8$.
\end{exercise}
